% Main figures -- controlled v117 scientific revision

\clearpage

\begin{figure}[p]
\centering
\includegraphics[width=\textwidth]{Fig1.png}
\caption{\textbf{Inventory adjudication and evidence-state contraction.} \textbf{a}, Reconciliation of the raw HydroLAKES scope with mapped reservoirs. The reference removes 108 authoritative and 645 calibrated high-confidence duplicate matches and restores 1{,}245 records without duplicate evidence. \textbf{b}, Generation retained from the reference L0 through L1, core and conservative L2 intervals, and partial L3 intervals. Lower and upper endpoints are evidence policies, not confidence limits. Complete published-criterion L2 and complete deployable L3 remain unavailable.}
\label{fig:framework}
\end{figure}

\begin{figure}[p]
\centering
\includegraphics[width=\textwidth]{Fig2.png}
\caption{\textbf{Geographic distribution of the reference and screened FPV resource.} \textbf{a}, Water-body-level L0 generation aggregated for global visualization. \textbf{b}, Continental L0 and L1 generation. \textbf{c}, Leading countries under L0, L1 and the core L2 lower endpoint. Country and continent assignments are complete for all 199{,}976 records; the map represents the study scope rather than all inland waters.}
\label{fig:geography}
\end{figure}

\begin{figure}[p]
\centering
\includegraphics[width=\textwidth]{Fig3.png}
\caption{\textbf{Model and resource-assumption evaluation.} \textbf{a}, Annual specific yield from the present monthly screen against the published Woolway per-kW output for 134{,}703 matched lakes. \textbf{b}, L0 and L1 generation across coverage scenarios. \textbf{c}, L1 generation under alternative ice-evidence policies. \textbf{d}, L0 response to wind treatments on the common legacy climate scope. The yield panel is a cross-model benchmark, not operating-plant validation.}
\label{fig:validation}
\end{figure}

\begin{figure}[p]
\centering
\includegraphics[width=\textwidth]{Fig4.png}
\caption{\textbf{Public-evidence L2 intervals.} \textbf{a}, Sequential contraction after core and conservative conservation and population screens. \textbf{b}, Generation associated with known no-dry-up, known observed-dry-up and unknown public records. Lower endpoints retain the first state; upper endpoints also retain the unknown state. \textbf{c}, Sensitivity of the L2 endpoints to three dry-up criteria applied to the 1991--2018 public GLEV record. The unresolved 2019--2020 gap is common to every criterion.}
\label{fig:l2_evidence}
\end{figure}

\begin{figure}[p]
\centering
\includegraphics[width=\textwidth]{Fig5.png}
\caption{\textbf{Uncertainty in partial infrastructure-proximity screening.} \textbf{a}, Core and conservative L2 and partial-L3 intervals. \textbf{b}, Core 2$\times$2 factorial separating dry-up evidence policy from road and power-line mapping policy. \textbf{c}, Generation retained with the GRanD hydropower-use proxy and with mapped lines alone. Power-line proximity and hydropower association do not establish voltage, substation access or residual interconnection capacity.}
\label{fig:accessibility}
\end{figure}
